\documentclass[11pt]{article}

\usepackage[margin=1in]{geometry}
\usepackage{amsmath,amssymb,amsthm}
\usepackage{algorithm}
\usepackage{algpseudocode}
\usepackage{enumitem}
\usepackage{hyperref}
\usepackage{xcolor}
\usepackage{tikz}
\usepackage{booktabs}

\hypersetup{colorlinks=true, linkcolor=blue, urlcolor=blue, citecolor=blue}

\newtheorem{theorem}{Theorem}[section]
\newtheorem{lemma}[theorem]{Lemma}
\newtheorem{proposition}[theorem]{Proposition}
\newtheorem{corollary}[theorem]{Corollary}
\newtheorem{definition}[theorem]{Definition}
\newtheorem{example}[theorem]{Example}

\title{TOAE209/TOAH209 -- Design and Analysis of Algorithms\\[4pt]
\large Week 13: PSPACE and Extending Tractability}
\author{Foreign Trade University}
\date{}

\begin{document}
\maketitle
\tableofcontents

\section{Introduction: Beyond NP}

The previous weeks built up a sharp picture of the world of computation: some
problems are solvable in polynomial time (the class \textbf{P}), many natural
problems are \textbf{NP-complete} and (conjecturally) admit no polynomial-time
algorithm, and reductions let us relate one hard problem to another. This week
we push the picture in two directions.

\begin{itemize}
    \item \textbf{Upward}, to a larger complexity class: \textbf{PSPACE}, the
    problems solvable using a polynomial amount of \emph{memory} (regardless of
    time). PSPACE contains all of NP, and its canonical complete problem --
    deciding a \emph{quantified} boolean formula (QSAT/QBF) -- captures the
    flavor of \emph{games}, where two players alternate moves. Game playing is
    naturally PSPACE-shaped: ``does there \emph{exist} a move for me such that
    for \emph{all} replies of my opponent, there exists a move for me, \dots,
    I win?''

    \item \textbf{Sideways}, by \emph{extending tractability}: even when a
    problem is NP-hard in general, we can often solve the instances we actually
    care about, \emph{provably} and \emph{efficiently}, by exploiting (a) a
    small \emph{parameter}, (b) special \emph{structure} in the input, or (c)
    clever \emph{pruning} of an exponential search. These are the themes of
    fixed-parameter tractability (FPT), dynamic programming on special graphs,
    and branch-and-bound.
\end{itemize}

This follows Kleinberg \& Tardos, Chapters 9 (PSPACE) and 10 (extending the
limits of tractability). The unifying message: \emph{NP-hardness is a
worst-case statement}. It does not forbid us from solving real instances; it
just tells us where to look for the leverage that makes a particular family of
instances tractable.

\section{PSPACE and Quantified Boolean Formulas}

\subsection{The class PSPACE}

\begin{definition}[PSPACE]
A decision problem is in \textbf{PSPACE} if there is an algorithm that decides
it using an amount of \emph{working memory} (space) bounded by a polynomial in
the input size $n$. There is no restriction on running time (it may be
exponential), only on space.
\end{definition}

Two basic facts position PSPACE relative to the classes we already know.

\begin{theorem}
$\textbf{P} \subseteq \textbf{NP} \subseteq \textbf{PSPACE}$.
\end{theorem}

\begin{proof}[Proof sketch]
$\textbf{P} \subseteq \textbf{NP}$ is immediate. For $\textbf{NP} \subseteq
\textbf{PSPACE}$: an NP problem can be decided by trying every candidate
certificate in turn and running the (polynomial-time) verifier on each. We can
enumerate certificates one at a time, reusing the same polynomial-size memory
for each verification, discarding it before the next -- so total \emph{space}
stays polynomial even though we examine exponentially many certificates over
exponential \emph{time}.
\end{proof}

The reuse-of-space idea is the heart of PSPACE: \emph{time can blow up
exponentially, but space is recycled}. It is widely believed (but unproven)
that $\textbf{NP} \neq \textbf{PSPACE}$, so PSPACE is plausibly strictly larger
than NP.

\subsection{Quantified Boolean Formulas (QBF / QSAT)}

Recall that SAT asks whether \emph{there exists} a satisfying assignment:
$\exists x_1 \exists x_2 \cdots \exists x_n\; \Phi(x_1,\dots,x_n)$. QSAT
generalizes this by allowing \emph{alternating} quantifiers.

\begin{definition}[QSAT / TQBF]
A \emph{quantified boolean formula} has the form
\[
Q_1 x_1\; Q_2 x_2 \cdots Q_n x_n\; \Phi(x_1,\dots,x_n),
\]
where each $Q_i \in \{\exists, \forall\}$ and $\Phi$ is a boolean formula
(often in CNF). The \textbf{QSAT} (also \textbf{TQBF}, ``true QBF'') problem is
to decide whether the fully-quantified formula evaluates to \texttt{true}.
\end{definition}

A fully-quantified formula has a definite truth value (no free variables). It
is evaluated by the obvious recursion: $\exists x_i\,\Psi$ is true iff $\Psi$
is true for $x_i=\texttt{false}$ \emph{or} for $x_i=\texttt{true}$;
$\forall x_i\,\Psi$ is true iff $\Psi$ is true for \emph{both} values. Problem
1 asks you to implement exactly this recursive evaluator.

\begin{theorem}[Stockmeyer--Meyer]
QSAT is PSPACE-complete.
\end{theorem}

\begin{proof}[Why it is in PSPACE]
The recursive evaluator above uses a recursion stack of depth $n$ (one frame
per quantifier), each frame storing $O(1)$ extra information (the current value
being tried for one variable). At a leaf we evaluate $\Phi$ on a complete
assignment in polynomial space. Although the recursion explores a tree of up to
$2^n$ leaves (exponential \emph{time}), at any instant only one root-to-leaf
path is materialized, so the \emph{space} is $O(n + |\Phi|)$ -- polynomial.
Completeness (that every PSPACE problem reduces to QSAT) is shown by encoding
a polynomial-space computation as a quantified formula; we omit the
construction (Kleinberg \& Tardos, Section 9.4).
\end{proof}

The alternation $\exists\forall\exists\forall\cdots$ is exactly the logical
structure of a two-player game, which we turn to next.

\section{Games as PSPACE Problems}

\subsection{Game trees and minimax}

Consider a finite, two-player, zero-sum, perfect-information game (tic-tac-toe,
Nim, chess on a bounded board, \dots). Players alternate moves; one player
(\textbf{Max}) wants to maximize the final payoff and the other (\textbf{Min})
wants to minimize it. The set of all play sequences forms a \textbf{game tree}:
internal nodes are positions (alternating between ``Max to move'' and ``Min to
move''), edges are moves, and leaves carry payoffs.

\begin{definition}[Minimax value]
The \emph{minimax value} of a position is defined recursively: a leaf's value is
its payoff; a Max node's value is the \emph{maximum} of its children's values; a
Min node's value is the \emph{minimum} of its children's values. The minimax
value of the root is the payoff achievable under optimal play by both sides.
\end{definition}

\begin{algorithm}[h]
\caption{Minimax (value for the player to move)}
\begin{algorithmic}[1]
\Function{Minimax}{$node$, $maximizing$}
    \If{$node$ is a leaf} \Return payoff$(node)$ \EndIf
    \If{$maximizing$}
        \State \Return $\max_{c \,\in\, \text{children}(node)} \textsc{Minimax}(c, \text{false})$
    \Else
        \State \Return $\min_{c \,\in\, \text{children}(node)} \textsc{Minimax}(c, \text{true})$
    \EndIf
\EndFunction
\end{algorithmic}
\end{algorithm}

The recursion mirrors a QBF exactly: ``does Max have a move ($\exists$) so that
for all Min replies ($\forall$) Max still has a move ($\exists$), \dots,
ending in a win?'' This is why deciding the winner of a generalized board game
is typically PSPACE-complete. Problem 2 asks you to implement minimax on a
nested-list game tree.

\subsection{Alpha-beta pruning}

Plain minimax explores the entire tree. \textbf{Alpha-beta pruning} computes
the \emph{same} minimax value while skipping subtrees that cannot affect the
result. It carries two bounds down the recursion:
\begin{itemize}
    \item $\alpha$ = the best (largest) value Max can already guarantee along
    the current path;
    \item $\beta$ = the best (smallest) value Min can already guarantee.
\end{itemize}
Whenever $\alpha \ge \beta$, the current node's remaining children are
irrelevant -- the opponent already has a better option elsewhere -- so we
\emph{prune}.

\begin{algorithm}[h]
\caption{Alpha-Beta Pruning}
\begin{algorithmic}[1]
\Function{AlphaBeta}{$node$, $maximizing$, $\alpha$, $\beta$}
    \If{$node$ is a leaf} \Return payoff$(node)$ \EndIf
    \If{$maximizing$}
        \State $v \gets -\infty$
        \For{each child $c$ of $node$}
            \State $v \gets \max(v, \textsc{AlphaBeta}(c, \text{false}, \alpha, \beta))$
            \State $\alpha \gets \max(\alpha, v)$
            \If{$\alpha \ge \beta$} \textbf{break} \Comment{$\beta$ cut-off} \EndIf
        \EndFor
        \State \Return $v$
    \Else
        \State $v \gets +\infty$
        \For{each child $c$ of $node$}
            \State $v \gets \min(v, \textsc{AlphaBeta}(c, \text{true}, \alpha, \beta))$
            \State $\beta \gets \min(\beta, v)$
            \If{$\beta \le \alpha$} \textbf{break} \Comment{$\alpha$ cut-off} \EndIf
        \EndFor
        \State \Return $v$
    \EndIf
\EndFunction
\end{algorithmic}
\end{algorithm}

\begin{theorem}
Alpha-beta pruning returns the same value as plain minimax.
\end{theorem}

\begin{proof}[Proof idea]
A subtree is pruned only when its parent already has a sibling value forcing
$\alpha \ge \beta$. In that situation, the value of the pruned subtree could
only move the node's $\max$/$\min$ in a direction the ancestor on the other side
would never choose, so it cannot change any value on the path to the root.
Hence the root value is unaffected by the pruning.
\end{proof}

With an ideal move ordering, alpha-beta examines only $O(b^{d/2})$ leaves
instead of $O(b^d)$ -- effectively doubling the search depth reachable in the
same time. Problem 3 asks you to implement alpha-beta and verify both that it
returns the same value as minimax and that it visits strictly fewer leaves.

\section{Extending Tractability I: Fixed-Parameter Tractability}

NP-hardness is a worst-case statement over \emph{all} inputs of size $n$. Often
the inputs we care about have a small \emph{parameter} $k$ -- the size of the
solution, the treewidth, the number of clusters -- even when $n$ is large. FPT
isolates the exponential blow-up in $k$ alone.

\begin{definition}[Fixed-parameter tractable]
A parameterized problem (with input size $n$ and parameter $k$) is
\textbf{fixed-parameter tractable (FPT)} if it can be solved in time
$f(k)\cdot \mathrm{poly}(n)$ for some function $f$ depending only on $k$. The
exponential cost (if any) is confined to $f(k)$; the dependence on $n$ stays
polynomial.
\end{definition}

The point: if $k$ is small (say $\le 20$) then $f(k)$ -- even $2^k$ -- is a
fixed constant, and the algorithm is genuinely efficient on large $n$, in stark
contrast to a $2^n$ brute force.

\subsection{Bounded search trees: Vertex Cover in $O(2^k\, n)$}

\begin{definition}[Vertex Cover]
A \emph{vertex cover} of a graph $G=(V,E)$ is a set $C \subseteq V$ such that
every edge has at least one endpoint in $C$. The decision problem: given $G$
and an integer $k$, is there a vertex cover with $|C| \le k$?
\end{definition}

Vertex Cover is NP-complete, but parameterized by $k$ it is FPT, via a beautiful
\textbf{bounded search tree}. The key observation:

\begin{lemma}[Branching rule]
For any edge $(u,v)$, every vertex cover must contain $u$ or $v$ (or both).
\end{lemma}

So we pick any uncovered edge $(u,v)$ and \emph{branch}: either $u \in C$ or
$v \in C$. Each branch fixes one vertex and decrements the budget $k$ by one,
giving a search tree of depth $\le k$ and branching factor $2$.

\begin{algorithm}[h]
\caption{FPT Vertex Cover (bounded search tree)}
\begin{algorithmic}[1]
\Function{VC}{$E$, $k$, $C$}
    \If{$E = \emptyset$} \Return $C$ \Comment{all edges covered} \EndIf
    \If{$k = 0$} \Return \textsc{None} \Comment{budget exhausted, edges remain} \EndIf
    \State pick any edge $(u,v) \in E$
    \For{$w \in \{u, v\}$}
        \State $E' \gets$ edges of $E$ not incident to $w$
        \State $R \gets \textsc{VC}(E', k-1, C \cup \{w\})$
        \If{$R \neq \textsc{None}$} \Return $R$ \EndIf
    \EndFor
    \State \Return \textsc{None}
\EndFunction
\end{algorithmic}
\end{algorithm}

\begin{theorem}
\label{thm:vc-fpt}
The bounded-search-tree algorithm decides whether a vertex cover of size $\le k$
exists in time $O(2^k \cdot (n+m))$, and returns such a cover if one exists.
\end{theorem}

\begin{proof}
\textbf{Correctness.} The recursion explores all ways of choosing, for each
edge encountered, which endpoint enters the cover. By the branching lemma, every
vertex cover is consistent with at least one such sequence of choices, so if a
cover of size $\le k$ exists, some leaf of the search tree at depth $\le k$
finds it; conversely every returned set covers all edges (we only return when
$E=\emptyset$) and has size $\le k$ (each recursive level adds one vertex and
decrements $k$). \textbf{Running time.} The recursion tree has depth $\le k$ and
branching factor $2$, hence $\le 2^k$ leaves and $O(2^k)$ nodes; each node does
$O(n+m)$ work to filter edges. Total $O(2^k(n+m))$.
\end{proof}

Note the contrast with the $\binom{n}{k}$ brute force over all $k$-subsets: the
search tree depends \emph{exponentially only on $k$}. Problems 4--5 ask you to
implement this decider and verify it against a brute-force minimum vertex cover
on small graphs.

\subsection{Kernelization}

A second pillar of FPT is \textbf{kernelization}: in polynomial time, shrink
the instance to an equivalent ``kernel'' whose size depends only on $k$, then
run any (even brute-force) algorithm on the small kernel.

For Vertex Cover, two reduction rules suffice for a first kernel:

\begin{itemize}
    \item \textbf{High-degree rule.} If a vertex $v$ has degree $> k$, then $v$
    must be in every cover of size $\le k$. (If $v$ were \emph{not} in $C$, then
    all $> k$ of its incident edges would need their \emph{other} endpoints in
    $C$, forcing $|C| > k$.) So put $v$ in the cover, delete $v$ and its edges,
    and decrement $k$.
    \item \textbf{Isolated-vertex rule.} A degree-$0$ vertex covers no edge and
    is never needed; remove it.
\end{itemize}

\begin{theorem}[Kernel size]
After applying the two rules exhaustively, a yes-instance has at most $k^2$
edges and at most $k^2 + k$ non-isolated vertices.
\end{theorem}

\begin{proof}
Once no high-degree vertex remains, every vertex has degree $\le k$. A cover of
size $\le k$ touches every edge, and each of its $\le k$ vertices covers
$\le k$ edges, so there are at most $k \cdot k = k^2$ edges. Each such edge has
$2$ endpoints, so at most $2k^2$ vertex-incidences and certainly
$O(k^2)$ non-isolated vertices remain.
\end{proof}

Thus the original $n$ drops out entirely: brute force on the kernel costs only a
function of $k$. Problem 6 asks you to implement these two reduction rules and
return the reduced instance.

\section{Extending Tractability II: Exploiting Special Structure}

Even an NP-hard problem can become polynomial when restricted to inputs with
special structure. The cleanest example is solving a hard graph problem on a
\textbf{tree}.

\subsection{Independent Set on trees via dynamic programming}

\begin{definition}[Independent Set]
An \emph{independent set} of $G=(V,E)$ is a set $S \subseteq V$ with no edge
between any two members. Maximum Independent Set asks for the largest such $S$.
It is NP-hard on general graphs.
\end{definition}

On a \emph{tree}, however, a root-to-leaf DP solves it in linear time. Root the
tree anywhere and, for each vertex $v$, compute two quantities over the subtree
rooted at $v$:
\[
\mathrm{incl}[v] = 1 + \!\!\sum_{c \,\in\, \mathrm{children}(v)}\!\! \mathrm{excl}[c],
\qquad
\mathrm{excl}[v] = \!\!\sum_{c \,\in\, \mathrm{children}(v)}\!\! \max(\mathrm{incl}[c], \mathrm{excl}[c]).
\]
The first is the best subtree solution \emph{including} $v$ (so no child may be
included); the second is the best \emph{excluding} $v$ (each child is free to be
included or not).

\begin{algorithm}[h]
\caption{Maximum Independent Set on a Tree}
\begin{algorithmic}[1]
\State root the tree at any vertex; compute a post-order of the vertices
\For{each vertex $v$ in post-order (children before parents)}
    \State $\mathrm{incl}[v] \gets 1 + \sum_{c} \mathrm{excl}[c]$
    \State $\mathrm{excl}[v] \gets \sum_{c} \max(\mathrm{incl}[c], \mathrm{excl}[c])$
\EndFor
\State \Return $\max(\mathrm{incl}[\text{root}], \mathrm{excl}[\text{root}])$
\end{algorithmic}
\end{algorithm}

\begin{theorem}
The DP computes a maximum independent set of an $n$-vertex tree in $O(n)$ time.
\end{theorem}

\begin{proof}
\textbf{Correctness} is by induction on subtree size: $\mathrm{incl}[v]$ and
$\mathrm{excl}[v]$ correctly record the best independent set of $v$'s subtree
under the two cases for $v$, given (by induction) correct values for all
children; including $v$ forbids including any child (hence $\mathrm{excl}[c]$),
while excluding $v$ leaves each child free (hence $\max$). The root case
$\max(\mathrm{incl},\mathrm{excl})$ allows or disallows the root.
\textbf{Time}: each vertex is processed once, doing work proportional to its
number of children; summing over the tree gives $O(n)$.
\end{proof}

This is a template: many NP-hard problems become tractable on trees (and, more
generally, on graphs of bounded \emph{treewidth}) precisely because a tree
admits such a clean subtree decomposition. Problem 7 asks you to implement the
tree DP and verify it against a brute-force maximum independent set.

\section{Extending Tractability III: Branch and Bound}

Branch-and-bound keeps an exponential search structure but prunes aggressively
using \emph{bounds}: an optimistic estimate (upper bound for maximization,
lower bound for minimization) of the best completion of a partial solution. If
the bound cannot beat the best complete solution found so far (the
\emph{incumbent}), the whole subtree is discarded. The answer is always exact;
only the search is faster.

\subsection{Branch-and-bound for the 0/1 Knapsack}

\begin{definition}[0/1 Knapsack]
Given $n$ items with values $v_i$ and weights $w_i$, and capacity $W$, choose a
subset maximizing total value subject to total weight $\le W$.
\end{definition}

We branch on each item (include / exclude) and bound with the
\textbf{fractional relaxation}: sort items by value-to-weight ratio, greedily
take whole items, then a fraction of the next to fill the remaining capacity.
This is the optimal value of the LP relaxation, hence a valid \emph{upper
bound} on any integral completion.

\begin{algorithm}[h]
\caption{Branch-and-Bound 0/1 Knapsack}
\begin{algorithmic}[1]
\State sort items by $v_i / w_i$ descending; $best \gets 0$
\Function{Explore}{$i$, $curVal$, $curWt$}
    \State $best \gets \max(best, curVal)$
    \If{$i = n$} \Return \EndIf
    \If{\textsc{FractionalBound}$(i, curVal, curWt) \le best$} \Return \Comment{prune} \EndIf
    \If{$curWt + w_i \le W$} \textsc{Explore}$(i{+}1, curVal+v_i, curWt+w_i)$ \EndIf
    \State \textsc{Explore}$(i{+}1, curVal, curWt)$
\EndFunction
\end{algorithmic}
\end{algorithm}

The optimum found must equal the value from the standard $O(nW)$ DP; Problem 8
asks you to verify this on random instances.

\subsection{Branch-and-bound for the TSP}

\begin{definition}[Travelling Salesman Problem]
Given a symmetric distance matrix on $n$ cities, find the shortest tour (cyclic
permutation) visiting every city exactly once and returning to the start.
\end{definition}

We branch on the order in which cities are appended to a partial tour and prune
with a \emph{lower bound} on the cost to complete it -- for instance, the cost
so far plus, for the current city and every unvisited city, the cheapest edge
that could still leave it. If this admissible bound already exceeds the best
complete tour found, the partial tour is abandoned.

\begin{algorithm}[h]
\caption{Branch-and-Bound TSP}
\begin{algorithmic}[1]
\State $best \gets +\infty$; start the tour at city $0$
\Function{Explore}{$current$, $count$, $cost$}
    \If{$cost \ge best$} \Return \EndIf
    \If{$count = n$} \State $best \gets \min(best, cost + d(current, 0))$; \Return \EndIf
    \If{\textsc{LowerBound}$(current, cost) \ge best$} \Return \Comment{prune} \EndIf
    \For{each unvisited city $j$}
        \State mark $j$ visited; \textsc{Explore}$(j, count+1, cost + d(current, j))$; unmark $j$
    \EndFor
\EndFunction
\end{algorithmic}
\end{algorithm}

The optimum must match an exact baseline such as Held--Karp DP (in
$O(2^n n^2)$) or brute force; Problem 9 asks you to verify this on small $n$.

\section{A Polynomial Island: 2-SAT}

Not every restriction of a hard problem stays hard. \textbf{2-SAT} -- SAT where
every clause has exactly two literals -- is solvable in \emph{linear} time,
even though $3$-SAT is NP-complete. The trick is the \textbf{implication
graph}: a clause $(a \lor b)$ is equivalent to the two implications
$(\lnot a \Rightarrow b)$ and $(\lnot b \Rightarrow a)$. Build a directed graph
with two nodes per variable ($x$ and $\lnot x$) and these implication edges.

\begin{theorem}[Aspvall--Plass--Tarjan]
A 2-CNF formula is satisfiable iff no variable $x$ has both $x$ and $\lnot x$
in the same strongly connected component (SCC) of the implication graph. If
satisfiable, a satisfying assignment can be read off in linear time from the
reverse topological order of the SCCs.
\end{theorem}

\begin{proof}[Proof idea]
Implications are transitive, so within an SCC all literals must take the same
truth value. If $x$ and $\lnot x$ are forced equal, the formula is
contradictory. Otherwise, processing SCCs in reverse topological order and
setting each not-yet-assigned literal to \texttt{true} (and its negation
\texttt{false}) never violates an implication, because an implication
$u \Rightarrow w$ cannot point from a later (true) SCC to an earlier (false)
one.
\end{proof}

Computing SCCs via Kosaraju's or Tarjan's algorithm gives the whole procedure
in $O(n+m)$. Problem 10 asks you to implement this and return an assignment or
\texttt{None}.

\section{State-Space Search: Iterative Deepening}

Many PSPACE-flavored puzzles (sliding tiles, word ladders) reduce to finding a
path in an enormous, implicitly-defined state graph. Breadth-first search finds
a shortest path but needs memory proportional to the frontier -- often
infeasible. \textbf{Iterative-deepening DFS (IDDFS)} combines BFS's optimality
(shortest path in number of moves) with DFS's small memory footprint: run a
depth-limited DFS with limit $1, 2, 3, \dots$ until the goal is found.

\begin{algorithm}[h]
\caption{Iterative-Deepening DFS}
\begin{algorithmic}[1]
\For{$\ell = 1, 2, 3, \dots$}
    \State run depth-limited DFS from the start with depth limit $\ell$
    \If{the goal is reached} \Return the path \EndIf
\EndFor
\end{algorithmic}
\end{algorithm}

Although shallow levels are re-explored, the cost is dominated by the deepest
level (a geometric series), so the total work is within a constant factor of a
single DFS to the solution depth, while space stays $O(\ell)$. Problem 11 asks
you to solve a small sliding puzzle with IDDFS and return the solution path.

\section{Set Cover by FPT/Bitmask Search}

\begin{definition}[Set Cover]
Given a universe $U$ and a family of subsets $S_1,\dots,S_m \subseteq U$, find
the fewest subsets whose union is $U$.
\end{definition}

Set Cover is NP-hard, but parameterized by the solution size $k$ it admits a
bounded-search-tree algorithm: pick any still-uncovered element $e$; some chosen
set must contain $e$, so branch over the (often few) sets containing $e$,
decrementing the budget. Representing covered elements as a \emph{bitmask}
makes union and coverage tests $O(1)$ machine operations. Problem 12 asks you
to implement this search parameterized by solution size and verify it against
brute force.

\section{Summary and Looking Ahead}

This week extended our map of the computational landscape in two directions.

\begin{itemize}
    \item \textbf{PSPACE} captures polynomial-\emph{space} computation and
    contains NP. Its complete problem, \textbf{QSAT/QBF}, has alternating
    $\exists/\forall$ quantifiers -- the logic of two-player games. \textbf{Game
    trees} are evaluated by \textbf{minimax}, sped up (without changing the
    answer) by \textbf{alpha-beta pruning}.

    \item \textbf{Fixed-parameter tractability} confines exponential cost to a
    parameter $k$: \textbf{bounded search trees} give Vertex Cover in
    $O(2^k n)$ (Theorem~\ref{thm:vc-fpt}), and \textbf{kernelization} shrinks an
    instance to size depending only on $k$.

    \item \textbf{Special structure} makes hard problems easy: Independent Set
    is linear-time on \textbf{trees} via DP, and \textbf{2-SAT} is linear-time
    via the implication-graph/SCC method even though 3-SAT is NP-complete.

    \item \textbf{Branch-and-bound} solves Knapsack and TSP \emph{exactly} by
    pruning with relaxation bounds, and \textbf{iterative-deepening} searches
    huge state spaces with little memory.
\end{itemize}

The throughline: \emph{NP-hardness is a worst-case verdict, not a death
sentence.} Real instances often have a small parameter, special structure, or
bounds tight enough to prune the search dramatically.

Next week (\textbf{Approximation Algorithms and Local Search}) we relax a
different requirement: instead of guaranteeing an \emph{exact} optimum quickly,
we will design polynomial-time algorithms that guarantee a solution provably
\emph{close} to optimal (approximation ratios), and study local-search
heuristics that iteratively improve a solution toward a local optimum.

\end{document}
